\section{Changing the normal weights}\label{sec:statement}
Let $C$ be a smooth connected projective curve of genus $g\ge2$, and
let $L\oplus M$ be a balanced bundle, meaning
$LM=K_C$ and $\deg L=\deg M=g-1$. Write $H$ for the rational
Chern-number-zero, virtual-dimension-zero summand of Pardon's semi-Fano universal cycle algebra \cite[Definition~1.1 and Theorem~1.2]{Pardon} and
$x_{h,1}\in H$ for its multiplicity-one scalar generator in genus $h$.
For two nonzero integral fiber weights with ratio $z$, let
$X_{g,m}(z)$ denote the normalized class of the multiple zero section
of covering degree $m$. Section~\ref{sec:weights} constructs its
Laurent polynomial in $z$. We use the logarithm for disjoint union
and write
\[
 P_{m,j}=[z^j q^m]\log\Big(1+\sum_{l>0}X_{g,l}(z)q^l\Big).
\]
The core genus $g$ is fixed in this notation. Let $N_{k,m,j}$ be the
coefficient of $z^j$ in the connected genus-$k$, degree-$m$ local
Gromov--Witten residue of $L\oplus M$, with the two independent fiber
weights as in \cite{BP}. Set $n_s=(g-1)\binom{s}{2}$ and
$w=(2\sin(u/2))^2$.

\begin{theorem}\label{thm:main}
Let $d$ be the largest proper divisor of $m\ge2$. If
$n_d<|j|\le n_m$, then
\[
 P_{m,j}=\sum_{h\ge0}a_{h,j}x_{h,1},\qquad
 \sum_{h\ge0}a_{h,j}w^{h-1}
       =\sum_{k\ge0}N_{k,m,j}u^{2k-2}.
\]
The left sum has finite support. Its coefficients vanish for
$h<1+m(g-1)+|j|$. For each genus bound $K$, the construction in
Theorem~\ref{thm:separate} gives $a_{h,j}$ for $h\le K$ as the
signed number of constrained embedded curves of genus $h$ divided by
$\prod_{i=2}^m i^{2(i-1)(g-1)}$.
\end{theorem}

For a prime $p$ one has $d=1$ and $n_d=0$. For the degree-$p$ logarithm $P_p(z)=\sum_jP_{p,j}z^j$,
symmetry gives
\[
        P_p(1)-P_p(-1)=4\sum_{\substack{j>0\\j\ {\rm odd}}}P_{p,j}.
\]
This determines the connected scalar--antidiagonal discrepancy at $p$.
At composite $m$ the boundary of the stated
range has a geometric meaning. If $d>1$, a linear system of spectral
polynomials at $j=n_d$ necessarily meets the locus of powers of
degree-$d$ polynomials. Division of cycles then detects a nonzero
lower-degree primitive. We give the intersection calculation in
\S\ref{sec:weights}.

Let $X$ be a smooth projective complex threefold with
$K_X\simeq\mathcal O_X$ and $H^1(X,\mathcal O_X)=H^2(X,\mathcal O_X)=0$.
For an effective integral class $\beta$, let $U_{X,\beta}\in H$ be
the class of its compact cycle space. Write $F_h^{\rm sc}$ and
$F_h^{\rm ad}$ for the scalar and balanced antidiagonal universal
local series of a genus-$h$ core, both with constant term one.
The exponents $e_{\beta,h}$ and $c_{\beta,h}$ are defined by
\[
 \sum_\beta U_{X,\beta}Q^\beta
    =\prod_{\beta\ne0,h}F_h^{\rm sc}(Q^\beta)^{e_{\beta,h}}
    =\prod_{\beta\ne0,h}F_h^{\rm ad}(Q^\beta)^{c_{\beta,h}}
\]
in the effective integral class monoid completed by ample degree.
The $e_{\beta,h}$ are rational, and $c_{\beta,h}$ are the integral
elementary-cluster coefficients by \cite[Theorem~1.1]{WueProduct}.
If the image of $\alpha$ in the free
quotient of $H_2(X;\ZZ)$ is primitive, write
$D_{g,p}(w)$ for the difference of the scalar and antidiagonal connected
degree-$p$ local residues of a core of genus $g$. Theorem~\ref{thm:main}
gives the coefficient identity
\begin{equation}\label{eq:compact-prime}
 \sum_h(c_{p\alpha,h}-e_{p\alpha,h})w^{h-1}
    = \sum_{p\gamma=p\alpha}\sum_{g\ge2}c_{\gamma,g}D_{g,p}(w).
\end{equation}
Indeed the only covering degrees at $p\alpha$ are one and $p$, and at
every effective integral root $\gamma$ the degree-one generators of both
products agree. For core genera zero and one the two local classes
agree in all degrees. Finite support at each class makes~\eqref{eq:compact-prime}
a polynomial identity after taking the difference.

The first possible change is in genus $p+2$. The use of the rational
basis $x_{h,1}$ in this argument can also be made before assigning any
numerical invariant to a complex family: proper pullback to connected
maps is injective on the multiplicity-one span when all source genera
are used \cite[Sections 5--7]{WueProduct}. For integral denominators we
use the stronger result that the $a_{h,d}:=X_{h,d}(-1)$ freely generate
$H_{\ZZ}$, with $a_{h,1}=x_{h,1}$. This is the integral basis theorem
of \cite[Theorem~1.1]{WueGenerators}. In particular the intersection of
this integral algebra with the rational span of the $x_{h,1}$ is
their free integral span: they are distinct degree-one variables
in that polynomial presentation.

The proof of Theorem~\ref{thm:main} occupies
Sections~\ref{sec:weights}--\ref{sec:comparison}. The argument needs
a family of targets to give the same count in two different theories.
For a prescribed bound on source genus we vary the target and impose
the spectral coefficient conditions until the constrained connected
images are regular embedded curves. The cohomological insertion
counts their signs. The relative virtual class on maps also sees all
contracted components, and each embedded image of genus $h$ contributes
the entire degree-one series $w^{h-1}$. This construction uses an
algebraic relative obstruction theory on the chosen family and
proper genus detection on the cycle spaces.
